% ============================================================
% RIWAYAT JUDUL (tidak di-include ke main.tex — catatan kerja)
% Kriteria: klaim utama harus kebaca dari judul; hindari jargon RSA;
% JANGAN mengklaim lebih dari yang data 6-tipe dukung.
% ============================================================
%
% FINAL (dipilih 2026-08-18, setelah finding 10v2 + 11):
%   "Readable, Faithful, Used: Three Dissociable Properties of
%    Demographic Identity in a Language Model"
%   Rasional: (a) persis memayungi 4 hasil — readable (probe > mulut),
%   faithful (peta L11 H16), used (lokus kausal 4/6 tipe), dan temuan
%   pemersatunya: ketiganya TIDAK saling meramalkan (disosiasi ganda);
%   (b) tidak ada klaim kuantitatif yang bisa terbantah oleh robustness
%   check; (c) "dissociable" = istilah baku neuropsikologi utk pola
%   bukti yang persis kita punya (double dissociation L1 vs mid-stack).
%
% GUGUR — terbantah data sendiri:
% 1. "Faithfully Arranged but Rarely Used" (working title lama)
%    -> "rarely used" salah: 4/6 tipe punya lokus kausal signifikan.
% 2. "Used in Proportion to Fidelity" / variasi "tracks fidelity"
%    -> terbantah finding 11 (Spearman -0.26..-0.54; EDU counterexample).
% 4. "A Single Attention Head Encodes the Geometry of Human Opinion
%    Groups—but the Model Doesn't Use It"
%    -> dua-duanya overclaim: head-nya kepakai kok utk membaca (probe),
%    dan "doesn't use" salah utk 4/6 tipe.
%
% MASIH HIDUP sebagai cadangan (kalau reviewer minta yang lebih catchy):
% 3. "The Model Knows More Than It Says: A Fidelity Map of Demographic
%    Identity in LLMs Against Real Survey Ground Truth"
%    -> didukung hasil probe (22-30% lebih dekat), TAPI "knows" perlu
%    kualifikasi (pengetahuannya kasar: urutan kelompok per-soal ~nol,
%    kalah dari baseline buta-kelompok). Berisiko diserang pakai data
%    kita sendiri — makanya bukan pilihan utama.
